\documentclass[10pt,a4paper]{amsart}
\usepackage[margin=19mm]{geometry}
\usepackage{amsmath,amssymb,mathtools}
\usepackage[scr=rsfs,cal=euler]{mathalfa}
\usepackage{xfrac}
\usepackage[unicode,hidelinks,bookmarks=false]{hyperref}
\newcommand{\ie}{i.e.\ }
\newcommand{\cf}{cf.\ }
\newcommand{\resp}{respectively\ }
\newcommand{\Subsection}{\subsection}
\providecommand{\nobreakdash}{\nobreak\hbox{-}}
\setlength{\emergencystretch}{2em}
\multlinegap0pt
\usepackage{bm}
\usepackage{amssymb}
\usepackage{mathtools}
\newtheorem{ass}{Assumption}
\newtheorem{thm}{Theorem}
\newtheorem{prop}{Proposition}
\newcommand{\E}{\mathbb E}
\newcommand{\Om}{\Omega}
\newcommand{\Tr}{\operatorname{Tr}}
\newcommand{\Z}{\mathbb Z}
\newcommand{\supp}{\operatorname{supp}}
\title{Controlled atomic limit and analytic density of states in a strongly attractive random Kronig-Penney model}
\author{Masahiro Kaminaga}
\date{Source: arXiv:2508.12664v3 (CC BY 4.0)}
\begin{document}
\maketitle

\noindent\emph{Source and licence.} Controlled atomic limit and analytic density of states in a strongly attractive random Kronig-Penney model. Version arXiv:2508.12664v3 (CC BY 4.0), distributed by arXiv under the Creative Commons Attribution 4.0 International licence, http://creativecommons.org/licenses/by/4.0/, as recorded in the arXiv licence metadata for this exact version and shown on the abstract page https://arxiv.org/abs/2508.12664v3. Full licence notice: \texttt{licenses/arxiv-2508.12664-CC-BY-4.0.txt}.\par\smallskip

\noindent\textit{Editorial scope.} Counted unit: the article's prop prop:spectrum (source0.tex lines 825--837). The article's complete original proof of it is delivered with it (source0.tex lines 839--906, one \texttt{proof} environment of 68 lines). No mathematics is written, repaired or re-derived: the statement and the proof are the article's own lines, in the article's own notation, and no retained line is substituted or re-flowed.  Retained uncounted as the source-local context the proof needs: lines 308--311; lines 435--439; lines 531--534; lines 592--595; lines 657--660; lines 663--669; lines 732--743.  Nothing was omitted from any retained span, so the package carries no omission record.  No retained line contains a citation, so the wrapper carries no bibliography and no bibliography entry.  No retained line contains a figure environment, an image-inclusion command or drawing code, so no image is delivered and the package declares no asset dependency.  Editorial formatting only, in addition: an AMS \texttt{amsart} preamble that restates the article's own declarations and macros used by the retained lines, namely the environments \texttt{ass}, \texttt{thm}, \texttt{prop} on separate counters, together with the standard packages those lines need.  The retained environments print in this document as Ass 1, Theorem 1, Proposition 1 in order of appearance, whereas the article prints the same items with the numbering of its own full document.  The complete native source of the article as distributed in the arXiv e-print for this exact version is bundled unchanged as \texttt{sources/183428/source0.tex}.
\begin{equation}
m_+(z)=\E\Bigl[\Tr\bigl(\chi_C(H_\omega-z)^{-1}\chi_C\bigr)\Bigr].
\label{mplus}
\end{equation}

\begin{equation}
 E\in\rho(H_\omega)\quad\Longleftrightarrow\quad
 0\in\rho(A_\omega(\zeta)).
\label{spectral-correspondence}
\end{equation}

\begin{ass}[Local analytic density]\label{ass:density}
The restriction of $\nu$ to $J_0$ has a density $h$.  The function $h$
extends holomorphically to $\Om_\delta(J)$ and continuously to its closure.
\end{ass}

\begin{equation}
M:=\frac{2C_\nu}{\delta}.
\label{Mdef}
\end{equation}

\begin{equation}
s_*:=\sup_{\zeta\in\Om}s(\zeta).
\label{sstar}
\end{equation}

\begin{ass}[Quantitative strong-attraction condition]\label{ass:small}
With $M$ and $s_*$ defined in \eqref{Mdef} and \eqref{sstar}, assume
\begin{equation}
Ms_*<1
\label{smallness-1}
\end{equation}
\end{ass}

\begin{thm}[Real-analyticity of the DOS]\label{thm:main}
Assume Assumptions~\ref{ass:density} and \ref{ass:small}.  Let
\begin{equation}
I=z(J)
 =\left(-\frac{1}{4a^2},-\frac{1}{4b^2}\right).
\label{energy-I}
\end{equation}
Then the upper-half-plane Stieltjes transform $m_+(z)$ in \eqref{mplus}
admits a holomorphic continuation through $I$.  In particular, the density
of states $n(E)$ exists and is real-analytic on $I$.  The IDS $N(E)$ is also
real-analytic on $I$.
\end{thm}

\begin{prop}[A nontrivial spectral interval]\label{prop:spectrum}
Assume, in addition to the hypotheses of Theorem~\ref{thm:main}, that
\begin{equation}
J\subset\operatorname{int}(\supp\nu).
\label{J-support}
\end{equation}
Here $\operatorname{int}$ denotes interior and $\sigma(B)$ denotes the spectrum
of an operator $B$.  Then
$$
I=z(J)\subset\sigma(H_\omega)
$$
almost surely.
\end{prop}

\begin{proof}
We construct a Weyl sequence for the reduced operator.  This argument does
not require a pointwise inverse on the real axis.  Fix first
$v\in J$ and put
$$
r=e^{-1/(2v)}\in(0,1).
$$
At the real parameter $v$, the convolution operator $T(v)$ has Fourier
multiplier
\begin{equation}
\tau_v(\theta)
 =2v\sum_{k=1}^{\infty}r^k\cos(k\theta)
 =2v\frac{r\cos\theta-r^2}
 {1-2r\cos\theta+r^2}.
\label{T-symbol}
\end{equation}
Thus $\tau_v(\theta_v)=0$ for
$\theta_v=\arccos r$.  Define
\begin{equation}
 c_L(n)=\frac{1}{\sqrt{2L+1}}e^{in\theta_v}
 \boldsymbol 1_{[-L,L]}(n),
\label{cutoff-wave}
\end{equation}
where $\boldsymbol 1_{[-L,L]}(n)$ is the indicator of the integer interval
$\{n\in\Z:|n|\le L\}$.  We have
\begin{equation}
 \|T(v)c_L\|\longrightarrow0
 \quad\hbox{as }L\longrightarrow\infty.
\label{T-Weyl}
\end{equation}
Indeed, the convolution kernel of $T(v)$ belongs to $\ell^1(\Z)$.
To verify \eqref{T-Weyl}, truncate this kernel to a range $R$.  Away from a
boundary of width $R$, its action on \eqref{cutoff-wave} is multiplication by
the truncated Fourier multiplier.  Since the full multiplier vanishes at
$\theta_v$, the absolute value of the truncated multiplier is bounded by the
$\ell^1$ norm of the omitted tail.  The relative size of the boundary tends
to zero as $L\to\infty$.  Letting first $L\to\infty$ and then $R\to\infty$
proves \eqref{T-Weyl}.

Since $v\in\supp\nu$, for every $\epsilon>0$ and every $L$ a prescribed
block of $2L+1$ variables has a positive probability of satisfying
\begin{equation}
 \max_{|n-k|\le L}|V_\omega(n)-v|<\epsilon.
\label{good-block}
\end{equation}
Fix sequences $L_j\to\infty$ and $\epsilon_j\downarrow0$.  For each $j$,
consider infinitely many mutually disjoint blocks of length $2L_j+1$.
The corresponding events \eqref{good-block} are independent and all have the
same positive probability.  The Borel-Cantelli lemma shows that infinitely
many of them occur almost surely.  Taking the intersection of these
probability-one events over $j$, we may choose one successful block for each
$j$, with centers $k_j$ tending to infinity.  Translate
$c_{L_j}$ to the block centered at $k_j$ and denote the resulting unit vector
by $c_j$.  Translation invariance of $T(v)$ and \eqref{good-block} give
\begin{equation}
 \|A_\omega(v)c_j\|
 \le \epsilon_j+\|T(v)c_{L_j}\|\longrightarrow0.
\label{A-Weyl}
\end{equation}
Hence $0\in\sigma(A_\omega(v))$.  By
\eqref{spectral-correspondence}, $z(v)\in\sigma(H_\omega)$.

To remove the dependence of the null set on $v$, apply the preceding argument
to a countable dense subset $J_{\mathrm d}$ of $J$ and intersect the resulting
probability-one events.  Almost surely, the closed set $\sigma(H_\omega)$
contains $z(J_{\mathrm d})$.  Since $z$ is continuous on $J$,
it also contains its closure and therefore contains $I=z(J)$.
\end{proof}

\end{document}
